\documentclass{beamer}

\usepackage{graphicx}
\usepackage{mathtools}
\usepackage{amssymb}
\usepackage{hyperref}

\usetheme{default}
\setbeamertemplate{navigation symbols}{}

\title{An Introduction to \LaTeX}
\author{Peter de Waal\\with material from Lennart Herlaar}
\institute{Department of Information and Computing Sciences}
\date{OIA \LaTeX{} workshop}

\AtBeginSection[]
{
  \begin{frame}{Outline}
    \tableofcontents[currentsection]
  \end{frame}
}

\begin{document}

\begin{frame}
    \titlepage
\end{frame}

\begin{frame}{Outline}
    \tableofcontents
\end{frame}

\section{Introduction}

\subsection{What is \LaTeX{} and why?}

\begin{frame}{Why \LaTeX?}
    If you merely want to produce a passably good document --- something
    acceptable and basically readable but not really beautiful --- a simpler
    system will usually suffice.

    \medskip
    With \TeX{} the goal is to produce the finest quality.

    \medskip
    Donald E. Knuth
\end{frame}

\begin{frame}{\TeX{}}
    \begin{itemize}
        \item Program for typesetting technical text
              \begin{itemize}
                  \item Specialised in formulas
                  \item Professional typographic design
                  \item Freely available
                  \item Stable
              \end{itemize}
        \item File format based on ASCII
              \begin{itemize}
                  \item Easy to send by e-mail
                  \item Works on all systems
              \end{itemize}
        \item Easily extendable through macros
    \end{itemize}
\end{frame}

\begin{frame}{\LaTeX{}}
    \begin{itemize}
        \item Macro package for \TeX{}
              \begin{itemize}
                  \item Freely available
                  \item Stable
              \end{itemize}
        \item Logical structure
        \item Widely used in the exact sciences
              \begin{itemize}
                  \item Publishers require the use of their own style files with
                        macros for their own layout
              \end{itemize}
    \end{itemize}
\end{frame}

\begin{frame}{\LaTeX{} workflow}
    \begin{itemize}
        \item Create ASCII file(s) with \LaTeX{} commands
        \item Create a DVI file from this with \texttt{latex} (Device Independent)
        \item Turn the DVI file into a PostScript file with \texttt{dvips} or a
              PDF with \texttt{dvi2pdf}, or
        \item Create a PDF file directly with \texttt{pdflatex}
        \item Output can be viewed with
              \begin{itemize}
                  \item DVI: \texttt{xdvi}, \texttt{yap}
                  \item PS: Ghostscript
                  \item PDF: Acrobat Reader
              \end{itemize}
    \end{itemize}
\end{frame}

\begin{frame}{\LaTeX{} tools}
    \begin{itemize}
        \item There are several \TeX{} distributions
        \item Any ASCII editor is fine
        \item \LaTeX{} is not WYSIWYG
        \item There are tools to make working with \LaTeX{} easier:
              \begin{itemize}
                  \item WYSIWYG editor (LyX)
                  \item \LaTeX{} shells or IDEs (WinShell, Context, TeXnicCenter)
                  \item Editors with syntax highlighting (Emacs)
              \end{itemize}
    \end{itemize}
\end{frame}

\section{Using \LaTeX{}}

\subsection{The basics}

\begin{frame}{Basic \LaTeX{} document}
    Start with a document class such as \texttt{article}.
\end{frame}

\begin{frame}{Basic \LaTeX{} document}
    A minimal file also needs \texttt{\string\begin\{document\}} and
    \texttt{\string\end\{document\}}.
\end{frame}

\begin{frame}{Basic \LaTeX{} document}
    Add packages before the document starts, for example
    \texttt{\string\usepackage[english,dutch]\{babel\}}.
\end{frame}

\begin{frame}{Basic \LaTeX{} document}
    You can then start structuring the content with sections.
\end{frame}

\begin{frame}{Basic \LaTeX{} document}
    Typical sectioning commands are \texttt{\string\section},
    \texttt{\string\subsection}, and \texttt{\string\subsubsection}.
\end{frame}

\begin{frame}{Basic \LaTeX{} document}
    Table of contents entries are generated automatically from the sectioning
    commands.
\end{frame}

\begin{frame}{Basic \LaTeX{} document}
    \begin{itemize}
        \item \texttt{\string\section\{Introduction\}}
        \item \texttt{\string\section\{An interesting section\}}
        \item \texttt{\string\section\{Summary\}}
        \item \texttt{\string\tableofcontents} can be used to print the outline
    \end{itemize}
\end{frame}

\subsection{The basics of formatting}

\begin{frame}{Text formatting}
    \begin{itemize}
        \item \LaTeX{} automatically handles:
              \begin{itemize}
                  \item automatic numbering of chapters, sections, and so on
                  \item correct spacing between letters, words, and sentences
                  \item ligatures
                  \item word breaking (with \texttt{\string\selectlanguage\{english\}})
                  \item justification of paragraphs and pages
              \end{itemize}
    \end{itemize}
\end{frame}

\begin{frame}{Input}
    \begin{itemize}
        \item Forbidden characters include \# \$ \% \^{} \& \_ \{ \} and the backslash.
        \item Literal output is usually written with \texttt{\string\verb} or the
              \texttt{verbatim} environment.
        \item Ready-made strings:
              \begin{itemize}
                  \item \texttt{\string\today} today
                  \item \texttt{\string\TeX} \TeX{}
                  \item \texttt{\string\LaTeX} \LaTeX{}
              \end{itemize}
    \end{itemize}
\end{frame}

\begin{frame}{Special characters}
    \begin{itemize}
        \item Quotes:
              \begin{itemize}
                  \item ``click 'Start' to begin'' becomes typographic quotes in \LaTeX{}
              \end{itemize}
        \item Hyphens:
              \begin{itemize}
                  \item X-rated
                  \item page 2--24
                  \item yes --- or is it no
                  \item $\pi$ and $-\pi$
              \end{itemize}
        \item Ellipsis: \ldots
        \item Accents and special symbols: \"{o} is o-umlaut, \^{o} is o-circumflex,
              \={o} is o-macron, \AE{} is AE
    \end{itemize}
\end{frame}

\begin{frame}{Fonts}
    \begin{tabular}{ll}
        \texttt{\string\textrm\{...\}}     & roman         \\
        \texttt{\string\textsf\{...\}}     & sans serif    \\
        \texttt{\string\texttt\{...\}}     & typewriter    \\
        \texttt{\string\textmd\{...\}}     & medium        \\
        \texttt{\string\textbf\{...\}}     & bold          \\
        \texttt{\string\textup\{...\}}     & upright       \\
        \texttt{\string\textit\{...\}}     & italic        \\
        \texttt{\string\textsl\{...\}}     & slanted       \\
        \texttt{\string\textsc\{...\}}     & small caps    \\
        \texttt{\string\emph\{...\}}       & emphasized    \\
        \texttt{\string\textnormal\{...\}} & document font
    \end{tabular}
\end{frame}

\begin{frame}{Special use of fonts}
    \begin{itemize}
        \item Emphasis
              \begin{itemize}
                  \item \underline{text} text
                  \item \emph{emphasis} emphasis
              \end{itemize}
        \item Selection of fonts and sizes
              \begin{itemize}
                  \item \textit{You can also emphasize something if the text is italic}
                  \item \textsf{or in sans-serif style}
                  \item \texttt{or in typewriter style}
                  \item \textbf{And also if it is bold!}
              \end{itemize}
    \end{itemize}
\end{frame}

\begin{frame}{Font sizes}
    \begin{tabular}{ll}
        \texttt{\string\tiny}         & tiny            \\
        \texttt{\string\scriptsize}   & extremely small \\
        \texttt{\string\footnotesize} & very small      \\
        \texttt{\string\small}        & small           \\
        \texttt{\string\normalsize}   & normal          \\
        \texttt{\string\large}        & large           \\
        \texttt{\string\Large}        & larger          \\
        \texttt{\string\LARGE}        & even larger     \\
        \texttt{\string\huge}         & huge            \\
        \texttt{\string\Huge}         & largest
    \end{tabular}
\end{frame}

\section{Document structure}

\begin{frame}{Titles, chapters, sections, and so on}
    \begin{itemize}
        \item Title: \texttt{\string\title}, \texttt{\string\author}, and
              \texttt{\string\maketitle}
        \item Chapters, sections, and so on
              \begin{itemize}
                  \item In the \texttt{article} document class: \texttt{\string\section},
                        \texttt{\string\subsection}, \texttt{\string\subsubsection},
                        \texttt{\string\appendix}
                  \item In the \texttt{book} and \texttt{report} document classes:
                        \texttt{\string\chapter}, \texttt{\string\part}
              \end{itemize}
    \end{itemize}
\end{frame}

\begin{frame}[fragile]{References}
    \begin{block}{Input}
        Labels let you refer back to a section and its page number later in the
        document.
    \end{block}

    \medskip
    \begin{quote}
        2. Preparation\par
        And here comes the text of this paragraph.\ldots\par
        3. Solution\par
        In section 2 on page 10 we wrote \ldots
    \end{quote}
\end{frame}

\begin{frame}[fragile]{Bibliography}
    \begin{block}{Input}
        Citations are inserted with a citation command and a bibliography
        environment.
    \end{block}

    \begin{block}{Output}
        D. Knuth [1] mentioned this \ldots

        Bibliography

        1 Knuth, Donald E.: The TeXBook (1984), Addison-Wesley Publishing
        Company, Reading (USA). 480 p.
    \end{block}
\end{frame}

\section{Lists and tables}

\begin{frame}[fragile]{Lists: enumerate, itemize, description}
    \begin{enumerate}
        \item first item
        \item second item
    \end{enumerate}

    \begin{itemize}
        \item an item in itemize
        \item[-] and another one!
    \end{itemize}

    \begin{description}
        \item[high] a description
        \item[low] final example
    \end{description}
\end{frame}

\begin{frame}[fragile]{Lists: enumerate, itemize, description (nested)}
    \begin{enumerate}
        \item first item
              \begin{itemize}
                  \item a subitem
                  \item[-] and another one!
              \end{itemize}
        \item second item
              \begin{description}
                  \item[high] description
                  \item[low] and another one
              \end{description}
    \end{enumerate}
\end{frame}

\begin{frame}[fragile]{Tables}
    \begin{block}{Input}
        A \texttt{tabular} environment aligns numbers and labels into columns.
    \end{block}

    \medskip
    \begin{center}
        \begin{tabular}{|r|l|}
            \hline
            \multicolumn{2}{|c|}{Expenses} \\
            \hline \hline
            1000 & travel                  \\
            3700 & lodging                 \\
            800  & conference              \\
            \hline
            5500 & total                   \\
            \hline
        \end{tabular}
    \end{center}
\end{frame}

\section{Mathematics}

\begin{frame}{Math mode}
    \begin{itemize}
        \item \LaTeX{} has a special mode for working with mathematical formulas:
              \begin{itemize}
                  \item Inline in text: between \texttt{\string$} signs
                  \item Standalone unnumbered formula: between \texttt{\string\[}
                        and \texttt{\string\]} (or between
                        \texttt{\string\begin\{displaymath\}} and
                        \texttt{\string\end\{displaymath\}})
                  \item Standalone numbered formula: between
                        \texttt{\string\begin\{equation\}} and
                        \texttt{\string\end\{equation\}}
              \end{itemize}
    \end{itemize}
\end{frame}

\begin{frame}[fragile]{Examples}
    Pythagoras proved that $a^2 = b^2 + c^2$. Moreover,

    \begin{equation}
        \label{e:vb1}
        \lim_{n \to \infty} \sum_{k=1}^n
        \frac{1}{k^2} = \frac{\pi^2}{6}.
    \end{equation}

    In equation~(\ref{e:vb1}) we see \ldots
\end{frame}

\begin{frame}{Examples}
    Pythagoras proved that $a^2 = b^2 + c^2$. Moreover,

    \[
        \lim_{n \to \infty} \sum_{k=1}^{n} \frac{1}{k^2} = \frac{\pi^2}{6}.
    \]

    In equation~(1) we see \ldots
\end{frame}

\begin{frame}{More ways to typeset mathematical formulas}
    \[
        y = \begin{cases}
            a     & \text{if } d > c      \\
            b + x & \text{in the morning} \\
            \text{all day}
        \end{cases}
    \]

    \[
        X = \begin{pmatrix}
            x_{11} & x_{12} & \cdots \\
            x_{21} & x_{22} & \cdots \\
            \vdots & \vdots & \ddots
        \end{pmatrix}
    \]

    \[
        \operatorname{corr}(X, Y) =
        \frac{\sum_{i=1}^{n}(x_i - \bar{x})(y_i - \bar{y})}
        {\sqrt{\sum_{i=1}^{n}(x_i - \bar{x})^2
                \sum_{i=1}^{n}(y_i - \bar{y})^2}}
    \]
\end{frame}

\section{Extensions}

\begin{frame}{EPS or PDF figures}
    \begin{itemize}
        \item \texttt{\string\usepackage\{graphicx\}}
        \item \texttt{\string\includegraphics[width=0.4\string\textwidth]\{figure\}}
        \item Create EPS with CorelDraw, FreeHand, xfig, ...
        \item EPS with \texttt{latex} and \texttt{dvips}, PDF with \texttt{pdflatex}
    \end{itemize}
\end{frame}

\begin{frame}{PDF, JPG and PNG figures}
    \begin{itemize}
        \item \texttt{\string\usepackage\{pgf\}}
        \item \texttt{\string\pgfdeclareimage[height=3.0cm]\{figuretwo\}\{fancy\}}
        \item \texttt{\string\pgfuseimage\{figuretwo\}}
        \item PDF from EPS using \texttt{epstopdf}, \texttt{ps2pdf}
    \end{itemize}
\end{frame}

\begin{frame}[fragile]{Floating figures and tables}
    \begin{block}{Input}
        A figure environment can carry a caption and a label.
    \end{block}

    We see in Figure~\ref{f:p2} \ldots
\end{frame}

\begin{frame}{Floating figures and tables}
    \begin{figure}
        \centering
        \fbox{\rule{0pt}{2cm}\rule{4cm}{0pt}}
        \caption{Image of Figure 2}
        \label{f:p2}
    \end{figure}

    We see in Figure~\ref{f:p2} \ldots
\end{frame}

\begin{frame}{Other extensions}
    \begin{itemize}
        \item Hyperlinks in PDF
              \begin{itemize}
                  \item Package \texttt{hyperref}
                  \item Hyperlinks for all references in the text
                  \item Hyperlinks to external references
              \end{itemize}
        \item PostScript tricks
              \begin{itemize}
                  \item Package \texttt{pstricks}
              \end{itemize}
        \item Managing references in a database
              \begin{itemize}
                  \item BibTeX with style files
              \end{itemize}
        \item Music
              \begin{itemize}
                  \item MusiXTEX package (macros plus fonts)
              \end{itemize}
    \end{itemize}
\end{frame}

\begin{frame}{Presentations}
    \begin{itemize}
        \item Package \texttt{beamer}
              \begin{itemize}
                  \item Creates a presentation PDF and handouts in one file
                  \item Hyperlinks between sections and to external URLs
                  \item Overlays
              \end{itemize}
    \end{itemize}
\end{frame}

\section{References}

\begin{frame}{References}
    \begin{itemize}
        \item L. Lamport,
              \emph{LaTeX: A Document Preparation System (2nd Edition)},
              Addison-Wesley, 1994.
        \item F. Mittelbach, M. Goossens, J. Braams, D. Carlisle, C. Rowley,
              \emph{The LaTeX Companion (2nd Edition)}, Addison-Wesley, 2004.
        \item T. Oetiker, H. Partl, I. Hyna, E. Schlegl,
              \emph{The Not So Short Introduction to \LaTeX{} 2\ensuremath{\varepsilon}},
              \url{http://www.ctan.org/tex-archive/info/lshort/english}
        \item CTAN, the Comprehensive \TeX{} Archive Network,
              \url{http://www.ctan.org}
    \end{itemize}
\end{frame}

\begin{frame}{Summary}
    \begin{itemize}
        \item \LaTeX{} is a text processing program:
              \begin{itemize}
                  \item Versatile
                  \item Customizable
                  \item Extensible
              \end{itemize}
    \end{itemize}
\end{frame}

\begin{frame}{Getting started\ldots}
    \begin{itemize}
        \item Practice room BBL-458.
        \item Use the WinShell program (in the Start menu under Tools).
        \item Get the exercises from the OIA website and try them.
    \end{itemize}
\end{frame}

\end{document}